\section{Analytic foundations and common models}\label{sec:foundations}

We will compare actual divisors on common models and later prove
smoothness at general points of flipped surfaces. This section
records the analytic facts needed for these two purposes. All
discrepancy comparisons use the fixed b-divisor $\bM$ from
Theorem~\ref{thm:termination}.

\subsection{Places and projective resolutions}\label{subsec:places}

A \emph{place} is a global prime divisor on a normal proper
bimeromorphic model, identified with its strict transforms on common
higher models. Its centre is its reduced irreducible image. A place
is exceptional over a model when its centre has codimension at least
two. This description does not identify places solely through fields
of global meromorphic functions. Properness makes the centres compact
analytic subspaces. When checking a local singularity condition, we
also use divisors on proper modifications of analytic neighbourhoods;
we refer to these explicitly as \emph{local places}.

Every divisor in a pullback comparison is an actual divisor, with
the compatible canonical representatives prescribed in
Theorem~\ref{thm:termination}. Ordinary relative canonical
coefficients can be computed from local pluricanonical equations
after clearing a Cartier multiple; over a smooth space they are
the orders of the Jacobian. These computations localize. Indeed,
a ratio of representatives that is a unit off a codimension-two
subset of a normal base extends as a unit across that subset.
Thus agreement on the isomorphic opens gives the same local
discrepancy coefficients.

For a proper bimeromorphic morphism $r:V\to T$ of normal spaces,
the set of points with positive-dimensional fibres is closed
analytic. Off this subset the morphism is finite and
bimeromorphic, hence an isomorphism by normality. The subset has
codimension at least two: its inverse image is a proper analytic
subset of $V$, and the fibre-dimension theorem excludes a
codimension-one image with positive-dimensional fibres. Fibres
are connected by Stein factorization. A point where the morphism
is a local isomorphism is isolated in its fibre, so it cannot
belong to a connected positive-dimensional fibre. The exceptional
locus is therefore exactly the inverse image of this subset.

We use normalization, proper mapping, fibre dimension, Stein
factorization, and resolution and principalization in the complex
analytic category. The last two permit simultaneous resolution
of finitely many divisors and coherent ideals by projective
modifications; see
\cite[Theorems~1.6, 1.10, and~13.2(1)]{BierstoneMilman1997}.
Projective morphisms are preserved under base change and
restriction to closed analytic subspaces; finite morphisms are
projective. The projective morphisms used here can also be
composed: a sufficiently large twist of one relative ample
line bundle by the pullback of the next is relatively ample
for the composite, with compactness giving a uniform choice.

\begin{proposition}[Common models]\label{prop:common-models}
Any finite portion of the sequence in
Theorem~\ref{thm:termination} has a smooth common model dominating
its members and $X'$, projective over each member. It may also
dominate any prescribed projective modification of one of those
members. For a chosen member $T$, it can be arranged that the
exceptional locus over $T$ is divisorial and that its union with
the strict transform of $\Supp B_T$ has simple normal crossings.
The model carries $\bM$ as the pullback of $M'$.
\end{proposition}

\begin{proof}
Enlarge the finite portion to an initial segment starting at $X_0$.
For one step, take the reduced component of
$X_i\times_{Z_i}X_{i+1}$ containing the graph over the common
isomorphic open. It is projective over both models, and its
normalization remains so. Combine these graphs and $X'\to X_0$
by successive dominating components of fibre products, then
normalize and resolve. To include a prescribed projective
modification of a member, take one further fibre product over
that member before resolving. Base change, restriction, finite
normalization, and composition show that all resulting morphisms
to the members of the segment are projective.

Fix $T$. Choose a closed analytic subset of codimension at
least two outside which the map is an isomorphism, $T$ is
smooth, and the boundary components are smooth and disjoint;
include their singular loci and pairwise intersections in this
subset. Principalize the pullback of its ideal and resolve
the total divisor without changing this open set. The inverse
image of the chosen subset is now a divisor whose components
are exceptional over $T$; it contains the entire exceptional
locus. Resolving together with the strict boundary gives the
claimed normal-crossing condition. Pulling back $M'$ preserves
the fixed data.
\end{proof}

A further common resolution with a model carrying any prescribed
place allows comparison of its discrepancy. That carrying model
need not itself be projective, and the further resolution is only
required to be proper. Projectivity of the models in
Proposition~\ref{prop:common-models} will be used when applying
negativity and when constructing divisorial witnesses.

\subsection{Negativity for actual divisors}

For a projective morphism, relative nefness of a rational divisor
in the following lemma is the degree condition on all contracted
curves. This is distinct from absolute analytic nefness.

\begin{lemma}[Analytic negativity]\label{lem:negativity}
Let $r:V\to T$ be a projective bimeromorphic morphism of normal
complex spaces and let $F$ be a $\Q$-Cartier divisor on $V$.
If $-F$ has nonnegative degree on every contracted curve and
$r_*F$ is effective, then $F$ is effective.
\end{lemma}

This is \cite[Lemma~1.3 and Appendix~A]{Wang2021}, after
clearing a Cartier multiple. The relative degree criterion in
the projective case is explicit in \cite[Appendix~B]{Wang2021}.
Effectivity is that of the actual divisor, not merely its class.
The corresponding statement is
\cite[Lemma~2.3(i)]{OrdinaryKahler}. We apply it only when $r_*F=0$.

An effective $\Q$-Cartier divisor on a normal space pulls back
effectively: after multiplying, its local meromorphic equation
has no divisorial poles, hence is holomorphic by normality.
It also has nonnegative degree on every compact curve not
contained in its support, by restriction to the normalization
of that curve. We will use both facts.

\subsection{Ordinary terminal singularities in codimension two}

The following local statement will be applied on an open set
meeting a flipped surface. We include its surface proof, as in
\cite[Lemma~4.4]{OrdinaryKahler}. Ordinary terminality here means
that every local exceptional place has ordinary log discrepancy
strictly greater than one.

\begin{lemma}\label{lem:terminal-smooth}
A normal complex fourfold with a canonical rational line bundle
and ordinary terminal singularities is smooth in codimension two.
\end{lemma}

\begin{proof}
Suppose the singular locus has a codimension-two component.
Near its general smooth point, choose a holomorphic projection
to $\C^2$ submersive on that component. Take a log resolution
which is an isomorphism over the regular locus. Choose a general
fibre $S$ of the projection: its inverse image $\widetilde S$
is smooth, misses exceptional strata whose images have dimension
less than two, and meets all remaining strata transversely. These
choices follow by generic smoothness on the resolution and its
finitely many strata.
Write $\rho:\widetilde S\to S$ for the induced resolution. Terminal
singularities are klt, so the ambient space has rational singularities
by \cite[Theorem~3.12]{FujinoAnalytic2022}, and is Cohen--Macaulay
\cite[Remark~2.1]{Greb2011}. The fibre has the expected codimension
two. Its two defining parameters therefore form a regular sequence,
so $S$ is Cohen--Macaulay. It is regular away from isolated points
by the general choice of fibre. Serre's criterion makes $S$ normal.

For the regular sequence, the Koszul calculation of the dualizing
sheaf gives $\omega_S\simeq\omega_T\otimes\mathcal O_S$, where
$T$ denotes the ambient neighbourhood; the two slice equations
trivialize the normal determinant. This uses the local analytic
Ext description of the dualizing sheaf
\cite[pp.~3--4]{RSW2017}. Away from the isolated intersections with
the ambient singular locus, a local Cartier pluricanonical power
restricts accordingly, and the identification extends reflexively
over the normal surface. Applying adjunction also on the smooth
resolution therefore restricts its relative canonical equation to
\[
 A:=K_{\widetilde S}-\rho^*K_S=\sum_j a_jE_j,
 \qquad a_j>0.
\]
Here each exceptional curve is a transverse restriction of an
exceptional divisor on the resolution; terminality gives its
positive coefficient.
A normal surface germ with this property is smooth. Indeed, if
exceptional curves remain, negative definiteness
\cite[Section~4.8(e)]{Grauert1962} gives $A^2<0$,
so $A\cdot E_j<0$ for some $j$. The pulled-back canonical divisor
has degree zero on $E_j$, and therefore
$K_{\widetilde S}\cdot E_j<0$. Since $E_j^2<0$, adjunction gives
\[
 0\leq p_a(E_j)
   =1+\tfrac12(K_{\widetilde S}\cdot E_j+E_j^2)\leq0.
\]
Both negative intersection numbers are integers, so each equals
$-1$. Thus $E_j$ is a smooth rational $(-1)$-curve. Contract it
to a smooth point by \cite[Section~4.8(a)]{Grauert1962}. The map to
the normal surface germ factors
through this contraction, and pushforward of the relative
canonical equation leaves positive coefficients on all remaining
exceptional curves. Repetition ends with a finite bimeromorphic
map to the normal surface germ, which is an isomorphism.

But the transverse slice at the chosen singular fourfold
point has embedding dimension at least that of the fourfold
germ minus two, hence greater than two. It is singular, a
contradiction.
\end{proof}

In Proposition~\ref{prop:ordinary-terminal} we will verify ordinary
terminality on an open set meeting each flipped surface before
applying the lemma there.
